\scanpage{50}
Générateurs et relations pour les groupoïdes et groupes de Teichmüller (spéciaux et ordinaires, évt. toriques...)

Je m'attends à trouver trois types de présentations, à propriétés complémentaires -- mais toutes à significations arithmétiques :

a) Présentation par ``pa\uline{ntalons}'' ou ``\uline{triades}''. Elle revient à travailler sur les pts à l'infini ``les plus critiques'' sur les multiplicités modulaires, et semble (d'après l'exemple de $\unsure{\mathfrak T|\mathscr T}_{0\,4}^{+!}$) commode surtout pour exprimer des groupes du type $T_{g,\nu}^{+!}$ (ou $\mathbb T_{g;\mathrm I}^{+}$).

b) Présentation par ``boîte'' complète de pièces de Légo (voir des boîtes encore plus complètes) à l'aide des \inserted{groupes d'}\uline{automorphismes} \inserted{\uline{stricts}} de ces pièces, et certains ``chemins'' de déformation particuliers allant de certaines pièces à d'autres. Exemples frappants : $T_{0,4}^+$ (sans !), ou $T_{1,1}^+\simeq\mathrm{Sl}(2,\mathbb Z)$ engendré par $\rho,\sigma$ avec $\rho^{\deleted{\unsure{2}}3}=\sigma^2$, $\rho^{\unsure{6|\sigma}}(=\sigma^4)=1$.

c) Présentations \uline{par involutions}, qui sont \unsure{dans|des} $\mathfrak T_{g,\nu}^{-}$. Exemple typique $T_{1\,1}\simeq\mathrm{Gl}(2,\mathbb Z)$ \deleted{\unsure{eng}} (\deleted{\unsure{Agir|agir}}\unsure{ine|es} en Groupes cartographiques triangulés) engendré par trois réflexions $r_1,r_2,r_3$ avec $(r_1r_2)^{\deleted{\unsure{4}}2}=(r_2r_3)^{\deleted{\unsure{5}}3},\ (r_1r_2)^4(=(r_2r_3)^6)=1$.

% archive sheet number: 52
